\documentclass[10pt,a4paper]{amsart}
\usepackage[margin=19mm]{geometry}
\usepackage{amsmath,amssymb,mathtools}
\usepackage[scr=rsfs,cal=euler]{mathalfa}
\usepackage{xfrac}
\usepackage[unicode,hidelinks,bookmarks=false]{hyperref}
\newcommand{\ie}{i.e.\ }
\newcommand{\cf}{cf.\ }
\newcommand{\resp}{respectively\ }
\newcommand{\Subsection}{\subsection}
\providecommand{\nobreakdash}{\nobreak\hbox{-}}
\setlength{\emergencystretch}{2em}
\multlinegap0pt
\usepackage{bm}
\usepackage{bbm}
\usepackage{dsfont}
\usepackage{xspace}
\usepackage{cancel}
\usepackage{mathtools}
\usepackage{amsthm}
\makeatletter
\@ifundefined{thm}{\newtheorem{thm}{Theorem}}{}
\makeatother
\DeclareMathOperator{\Ewens}{Ewens}
\DeclareMathOperator{\Hoppe}{Hoppe}
\DeclareMathOperator{\Prob}{Prob}
\DeclareMathOperator{\Y}{\mathbb{Y}}
\title[A refinement of the Ewens sampling formula]{A refinement of the Ewens sampling formula}
\author{Eugene Strahov}
\date{Source: arXiv:2403.05077v3 (CC BY 4.0)}
\begin{document}
\maketitle

\noindent\emph{Source and licence.} A refinement of the Ewens sampling formula. Version arXiv:2403.05077v3 (CC BY 4.0), distributed under the Creative Commons Attribution 4.0 International licence, as recorded in the arXiv licence metadata for this exact version and shown on the abstract page https://arxiv.org/abs/2403.05077v3. Full rights notice: licenses/arxiv-2403.05077-CC-BY-4.0.txt.\par\smallskip

\noindent\textit{Editorial scope.} Counted unit: the Thm whose source label is \texttt{TheoremGeneralizedHoope} in the article (source0.tex lines 877--888, source label \texttt{TheoremGeneralizedHoope}), delivered together with the article's own complete original proof of it (source0.tex lines 889--1009, one \texttt{proof} environment of 121 lines). No mathematics is written, repaired or re-derived: statement and proof are the article's own text line for line. Required context retained uncounted: MultipleEwensMeasure (lines 487--497); EwensMeasure (lines 499--501). Every definition, display and label of those lines is retained unchanged. Omitted from this package and not redistributed: 1--486, 498--498, 502--876, 1010--2618. Nothing inside a retained excerpt is omitted or altered: every retained body is delivered exactly as the article's lines are written, so no omission receipt of any kind accompanies this package. Citations: the retained excerpts carry no citation command at all, so no bibliography entry of the article is transcribed and no reference is redistributed. Reference adaptation: none was needed and none was made; every reference inside the delivered unit resolves inside it, so no unresolved reference or citation is produced. Printed slips kept literal: no printed word, symbol, hypothesis or number of the counted statement or of its proof was altered. Layout and class: the article's own class is \texttt{amsart} with packages including epsfig, amsfonts, color, graphicx, amsmath, amssymb; the wrapper is the lane's pinned \texttt{amsart} editorial wrapper at 10pt on A4, which restates the theorem-like environments the retained text needs and adds the article's own macro definitions that the retained text needs (\texttt{\textbackslash Ewens}, \texttt{\textbackslash Hoppe}, \texttt{\textbackslash Prob}, \texttt{\textbackslash Y}), with extra wrapper packages (bm, bbm, dsfont, xspace, cancel, mathtools). The delivered numbering is the unit's own. Assets: no retained span contains a figure environment, a graphics-inclusion command, a PDF-page inclusion, a table or a TikZ picture; no image file is copied into this package, the package declares no asset dependency and no image file is redistributed with it. Licence: version arXiv:2403.05077v3 of this article is distributed under the Creative Commons Attribution 4.0 International licence, as recorded in the arXiv licence metadata for this exact version and shown on the abstract page https://arxiv.org/abs/2403.05077v3. A full rights notice is delivered at licenses/arxiv-2403.05077-CC-BY-4.0.txt. Characters: every retained excerpt is pure ASCII, so the delivered engine, XeLaTeX with its default Latin Modern text font, reports no missing character.
\begin{equation}\label{MultipleEwensMeasure}
\begin{split}
&M_{\theta_1,\ldots,\theta_k;n}^{\Ewens}\left(\Lambda_n^{(k)}\right)
=\frac{\left(\theta_1\right)_{|\lambda^{(1)}|}\ldots \left(\theta_k\right)_{|\lambda^{(k)}|
}}{\left(\theta_1+\ldots+\theta_k\right)_n}\\
&\times\frac{n!}{|\lambda^{(1)}|!\ldots |\lambda^{(k)}|!}\;
M_{\theta_1;|\lambda^{(1)}|}^{\Ewens}\left(\lambda^{(1)}\right)
\ldots M_{\theta_k;|\lambda^{(k)}|}^{\Ewens}\left(\lambda^{(k)}\right),\;\;
\Lambda_n^{(k)}=\left(\lambda^{(1)},\ldots,\lambda^{(k)}\right).
\end{split}
\end{equation}

\begin{equation}\label{EwensMeasure}
M_{\theta; n}^{\Ewens}\left(\lambda\right)=\frac{n!}{\prod_{j=1}^{n}j^{m_j(\lambda)}m_j(\lambda)!}\frac{\theta^{l(\lambda)}}{(\theta)_n},
\end{equation}

\begin{thm}\label{TheoremGeneralizedHoope}
Let $\lambda^{(1)}$, $\ldots$, $\lambda^{(k)}$ be fixed
Young diagrams such that
$$
\left|\lambda^{(1)}\right|+\ldots+\left|\lambda^{(k)}\right|=n.
$$
Then $\Lambda_n^{(k)}=\left(\lambda^{(1)},\ldots,\lambda^{(k)}\right)$ is a fixed multiple partition of $n$ into $k$ components. We have
\begin{equation}
\Prob\left(\mathfrak{L}_n^{(k)}=\Lambda_n^{(k)}\right)=M_{\theta_1,\ldots,\theta_k;n}^{\Ewens}\left(\Lambda_n^{(k)}\right),
\end{equation}
where $n=1,2,\ldots$, and the probability measure $M_{\theta_1,\ldots,\theta_k;n}^{\Ewens}$ is defined by equations (\ref{MultipleEwensMeasure}) and (\ref{EwensMeasure}).
\end{thm}

\begin{proof}
We proceed by induction. If $n=1$, the possible multiple partitions are
$$
((1),\varnothing,\ldots,\varnothing),\; (\varnothing, (1),\varnothing,\ldots,\varnothing),\ldots,
(\varnothing,\varnothing,\ldots,(1)).
$$
In particular, the multiple partition
$$
\left(\underset{l}{\underbrace{\varnothing,\ldots,\varnothing, (1)}},\varnothing,\ldots,\varnothing\right)
$$
has probability $\frac{\theta_l}{\theta_1+\ldots+\theta_k}$. On the other hand,
\begin{equation}
\begin{split}
&M^{\Ewens}_{\theta_1,\ldots,\theta_k; n=1}\left(\left(\underset{l}{\underbrace{\varnothing,\ldots,\varnothing, (1)}},\varnothing,\ldots,\varnothing\right)
\right)\\
&=\frac{(\theta_l)_1}{(\theta_1+\ldots+\theta_k)_1}
\frac{1!}{0!\ldots 1!\ldots 0!}
M^{\Ewens}_{\theta_1,|\lambda^{(1)}|=0}\left(\varnothing\right)
\ldots
M^{\Ewens}_{\theta_l,|\lambda^{(l)}|=1}\left((1)\right)
\ldots
M^{\Ewens}_{\theta_k,|\lambda^{(k)}|=0}\left(\varnothing\right)\\
&=\frac{\theta_l}{\theta_1+\ldots+\theta_k},
\end{split}
\end{equation}
so the result holds true.

Denote by $P_{n}^{\Hoppe}$ the probability measure describing the distribution of the random multiple
partition $\mathfrak{L}_n^{(k)}$. We have $\Prob\left(\mathfrak{L}_n^{(k)}=\Lambda_n^{(k)}\right)=P_n^{\Hoppe}(\Lambda_n^{(k)})$, and
\begin{equation}\label{P1}
P_n^{\Hoppe}\left(\Lambda_n^{(k)}\right)=\sum\limits_{\Lambda_{n-1}^{(k)}\in\Y_{n-1}^{(k)}}
P^{\Hoppe}_{n-1}\left(\Lambda_{n-1}^{(k)}\right)p\left(\Lambda_{n-1}^{(k)},\Lambda_n^{(k)}\right),
\end{equation}
where
$$
p\left(\Lambda_{n-1}^{(k)},\Lambda_n^{(k)}\right)
=\Prob\left(\mathfrak{L}_n^{(k)}=\Lambda_n^{(k)}|\mathfrak{L}_{n-1}^{(k)}=\Lambda_{n-1}^{(k)}\right)
$$
is the transition probability of the Markov process $\left\{\mathfrak{L}_j^{(k)}\right\}_{j=1}^{\infty}$. Let us find this transition probability explicitly.
We rewrite equation (\ref{P1}) as
\begin{equation}
\begin{split}
&P^{\Hoppe}_n\left(\left(\lambda^{(1)},\ldots,\lambda^{(k)}\right)\right)\\
&=\sum\limits_{\mu^{(1)}\nearrow\lambda^{(1)}}P_{n-1}^{\Hoppe}\left(\left(\mu^{(1)},\lambda^{(2)},\ldots,\lambda^{(k)}\right)\right)
p\left(\left(\mu^{(1)},\lambda^{(2)},\ldots,\lambda^{(k)}\right),\left(\lambda^{(1)},\lambda^{(2)},\ldots,\lambda^{(k)}\right)\right)\\
&+\sum\limits_{\mu^{(2)}\nearrow\lambda^{(2)}}P_{n-1}^{\Hoppe}\left(\left(\lambda^{(1)},\mu^{(2)},\ldots,\lambda^{(k)}\right)\right)
p\left(\left(\lambda^{(1)},\mu^{(2)},\ldots,\lambda^{(k)}\right),\left(\lambda^{(1)},\lambda^{(2)},\ldots,\lambda^{(k)}\right)\right)\\
&+\ldots\\
&+\sum\limits_{\mu^{(k)}\nearrow\lambda^{(k)}}P_{n-1}^{\Hoppe}\left(\left(\lambda^{(1)},\lambda^{(2)},\ldots,\mu^{(k)}\right)\right)
p\left(\left(\lambda^{(1)},\lambda^{(2)},\ldots,\mu^{(k)}\right),\left(\lambda^{(1)},\lambda^{(2)},\ldots,\lambda^{(k)}\right)\right),
\end{split}
\end{equation}
where $\mu^{(l)}\nearrow\lambda^{(l)}$ means that $\lambda^{(l)}$ is obtained from $\mu^{(l)}$ by adding one box.
In order to get a formula for the transition probability
$$
p\left(\left(\lambda^{(1)},\ldots,\mu^{(l)},\ldots,\lambda^{(k)}\right),\left(\lambda^{(1)},\ldots,\lambda^{(l)},\ldots,\lambda^{(k)}\right)\right)
$$
(where $\mu^{(l)}\nearrow\lambda^{(l)}$ ) we observe that $\lambda^{(l)}$ can be obtained from $\mu^{(l)}$ in two ways:\\
(a) A new box is attached to the bottom of $\mu^{(l)}$. If $\lambda^{(l)}
=\left(1^{m_1^{(l)}}2^{m_2^{(l)}}\ldots\right)$, then
$\mu^{(l)}=\left(1^{m_1^{(l)}-1}2^{m_2^{(l)}}\ldots\right)$.
This corresponds to the addition of an object of class $l$ with a new color.
In this case, the transition probability for the generalized Hoppe urn model is
$$
\frac{\theta_l}{\theta_1+\ldots+\theta_k+n-1}.
$$
(b) A row of length $j$ in $\mu^{(l)}$ becomes a row of length $j+1$ in $\lambda^{(l)}$, where $j$ takes values in the set
$\left\{1,2,\ldots,\lambda_1^{(l)}-1\right\}$. If
$$
\lambda^{(l)}=\left(1^{m_1^{(l)}}2^{m_2^{(l)}}...\right),
$$
then
$$
\mu^{(l)}=\left(1^{m_1^{(l)}}2^{m_2^{(l)}}\ldots (j-1)^{m_{j-1}^{(l)}}j^{m_j^{(l)}+1}(j+1)^{m_{j+1}^{(l)}-1}(j+2)^{m_{j+2}^{(l)}}\ldots\right).
$$
In this case, at time $n-1$ a colored (non-black) object of class $l$ is chosen from the existing $j$ objects of class $l$. The chosen object is added with an additional object of the same class, and of the same color. The corresponding transition probability is
$$
\frac{j\left(m_j^{(l)}+1\right)}{\theta_1+\ldots+\theta_k+n-1}.
$$
It follows that the distribution $P^{\Hoppe}_n$ of the generalized Hoppe urn satisfies the recursion
\begin{equation}
\begin{split}
&1=\sum\limits_{l=1}^k\frac{\theta_l}{\theta_1+\ldots+\theta_k+n-1}
\frac{P^{\Hoppe}_{n-1}\left(1^{m_1^{(1)}}2^{m_2^{(1)}}\ldots;1^{m_1^{(l)}-1}2^{m_2^{(l)}}\ldots;
1^{m_1^{(k)}}2^{m_2^{(k)}}\ldots\right)}{P^{\Hoppe}_{n}\left(1^{m_1^{(1)}}2^{m_2^{(1)}}\ldots;1^{m_1^{(l)}}2^{m_2^{(l)}}\ldots;
1^{m_1^{(k)}}2^{m_2^{(k)}}\ldots\right)}\\
&+\sum\limits_{l=1}^k\sum\limits_{j=1}^{\lambda_1^{(l)}-1}
\frac{j(m_j^{(l)}+1)P^{\Hoppe}_{n-1}\left(1^{m_1^{(1)}}2^{m_2^{(1)}}\ldots;\ldots;1^{m_1^{(l)}}\ldots j^{m_j^{(l)}+1}(j+1)^{m_{j+1}^{(l)}-1}\ldots;\ldots;
1^{m_1^{(k)}}2^{m_2^{(k)}}\ldots\right)}{(\theta_1+\ldots+\theta_k+n-1)P^{\Hoppe}_{n}\left(1^{m_1^{(1)}}2^{m_2^{(1)}}\ldots;\ldots;1^{m_1^{(l)}}2^{m_2^{(l)}}\ldots;
\ldots;1^{m_1^{(k)}}2^{m_2^{(k)}}\ldots\right)}.
\end{split}
\nonumber
\end{equation}
Let us consider the probability measure $M^{\Ewens}_{\theta_1,\ldots,\theta_k;n}$ defined by equations (\ref{MultipleEwensMeasure}) and (\ref{EwensMeasure}). It is enough to show that
$M^{\Ewens}_{\theta_1,\ldots,\theta_k;n}$ satisfies the same recursion as $P^{\Hoppe}_{n}$.
We have
\begin{equation}
\frac{M^{\Ewens}_{\theta_1,\ldots,\theta_k;n-1}\left(1^{m_1^{(1)}}2^{m_2^{(1)}}\ldots;1^{m_1^{(l)}-1}2^{m_2^{(l)}}\ldots;
1^{m_1^{(k)}}2^{m_2^{(k)}}\ldots\right)}{M^{\Ewens}_{\theta_1,\ldots,\theta_k;n}\left(1^{m_1^{(1)}}2^{m_2^{(1)}}\ldots;1^{m_1^{(l)}}2^{m_2^{(l)}}\ldots;
1^{m_1^{(k)}}2^{m_2^{(k)}}\ldots\right)}=\frac{m_1^{(l)}\left(\theta_1+\ldots+\theta_k+n-1\right)}{n\theta_l},
\end{equation}
and
\begin{equation}
\begin{split}
&\frac{M^{\Ewens}_{\theta_1,\ldots,\theta_k;n-1}\left(1^{m_1^{(1)}}2^{m_2^{(1)}}\ldots;\ldots;1^{m_1^{(l)}}\ldots j^{m_j^{(l)}+1}(j+1)^{m_{j+1}^{(l)}-1}\ldots;\ldots;
1^{m_1^{(k)}}2^{m_2^{(k)}}\ldots\right)}{M^{\Ewens}_{\theta_1,\ldots,\theta_k;n}\left(1^{m_1^{(1)}}2^{m_2^{(1)}}\ldots;\ldots;1^{m_1^{(l)}}2^{m_2^{(l)}}\ldots;
\ldots;1^{m_1^{(k)}}2^{m_2^{(k)}}\ldots\right)}\\
&=\frac{\theta_1+\ldots+\theta_k+n-1}{n}\;\frac{j+1}{j}
\frac{m_{j+1}^{(l)}}{m_j^{(l)}+1}.
\end{split}
\end{equation}
We see that $M^{\Ewens}_{\theta_1,\ldots,\theta_k;n}$ satisfies the same recurrence relation as $P^{\Hoppe}_{n}$ provided that the condition
\begin{equation}\label{condition1}
1=\sum\limits_{l=1}^{k}\frac{m_1^{(l)}}{n}+\sum\limits_{l=1}^k\sum\limits_{j=1}^{\lambda_1^{(l)}-1}\frac{(j+1)m_{j+1}^{(l)}}{n}.
\end{equation}
is satisfied. Since
$$
m_1^{(l)}+\sum\limits_{j=1}^{\lambda_1^{(l)}-1}(j+1)m_{j+1}^{(l)}=\sum\limits_{j=1}^{\lambda_1^{(l)}}jm_j^{(l)}=|\lambda^{(l)}|,
$$
and $|\lambda^{(1)}|+\ldots+|\lambda^{(k)}|=n$, condition (\ref{condition1}) holds true indeed.
\end{proof}

\end{document}
